\scanpage{187}
\mspage{132}
% archive sheet number: 186

reste en suspens\unsure{.|,} pour l'\unsure{instant}. (Je soupçonne que c'est vrai, mais il faudrait, pour le prouver, tenir une description des flèches des $\mathrm{HOT}_W(S)$ par des chemins de longueur bornée, disons \deleted{\unsure{par}} de longueur deux ou trois. C'est une question intéressante, mais que je réserve pour plus tard, si le besoin m'y contraint\ldots

Plus importante est la question de la validité de Der 3 et de ses variantes. (I p. 111-112) Le minimum qu'on demande d'un dérivateur $\mathbb D$, c'est que si $S'\longrightarrow S$ dans $\mathrm{Diag}$ est une Hot-équivalence, ce soit une $\mathbb D$-équivalence, i.e. \deleted{que} (dans le cas actuel) que le foncteur

\[
f^*:\mathrm{HOT}_W(S)\longrightarrow\mathrm{HOT}_W(S')
\]

soit ple fid. sur la sous-catégorie des objets \uline{constants} de $\mathrm{HOT}_W(S)$. Cette condition devient plus tangible quand on dispose déjà \inserted{soit} des foncteurs cohomologiques $H^{\unsure{\bullet|\circ}}_{\mathbb D}(S,-)$, $H^{\unsure{\cdot|\prime}}_{\mathbb D}(S',-)$, soit de l'homologie $H^{\mathbb D}_{\cdot}(S,-),H^{\mathbb D}_{\cdot}(S',-)$. En termes de cette dernière, elle signifie que

\[
\deleted{\unsure{f}}\quad p_!p^*(a)\unsure{\overset{\sim}{\text{---}}|\xrightarrow{\sim}}p'_!p'^*(a)\qquad\text{iso},
\]
